% ch23 — Discontinuous Galerkin Methods ★
\chapter{Discontinuous Galerkin Methods\starmark{}}
\label{ch:dg-methods}

Discontinuous Galerkin (DG) methods combine the high-order accuracy of spectral
methods with the geometric flexibility and shock-capturing ability of finite
volumes. This advanced chapter develops the DG framework as implemented in the
codebase and applies it to wave propagation and the CCZ4 system.

\section{Weak Form of Conservation Laws}
\label{sec:dg-weak}
\coderef{dg}{rs}

We derive the weak (variational) form of a hyperbolic conservation law on a
single element, introducing test functions and integration by parts to produce
the volume and surface integral terms.

\section{Legendre-Gauss-Lobatto Basis}
\label{sec:dg-lgl}

We construct the nodal DG basis using Lagrange polynomials on LGL nodes,
connect to the spectral-element mass and stiffness matrices, and exploit
the diagonal mass-lumping approximation.

\section{Rusanov (Local Lax-Friedrichs) Numerical Flux}
\label{sec:dg-rusanov}

We define the Rusanov numerical flux at element interfaces, estimate the
maximum wave speed, and discuss its role in providing upwind dissipation
for stability.

\section{\texorpdfstring{$hp$}{hp}-Adaptivity}
\label{sec:dg-hp}
\coderef{dg\_hp}{rs}

We describe the $hp$-refinement strategy that adapts both the mesh size $h$
and the polynomial order $p$ based on a smoothness indicator, concentrating
degrees of freedom where they are most needed.

\section{Mortar Elements}
\label{sec:dg-mortar}
\coderef{dg\_mortar}{rs}

We introduce mortar elements as the coupling mechanism between non-conforming
element interfaces that arise from $h$-refinement, ensuring conservation
across hanging nodes.

\section{Cubed-Sphere Grids\starmark{}}
\label{sec:dg-cubed-sphere}
\coderef{dg\_cubed\_sphere}{rs}

We describe the cubed-sphere coordinate mapping that avoids polar singularities
for problems on spherical domains, and show how it integrates with the DG
framework for wave-extraction spheres.

\paragraph{Key References.}
The DG formulation follows \cite{HesthavenWarburton2008,Kopriva2009};
$hp$-adaptivity follows \cite{Mavriplis1994}; mortar elements follow
\cite{Kopriva1996}; the cubed-sphere mapping draws on
\cite{Ronchi1996}; the CCZ4 application follows \cite{Dumbser2018}.
\coderef{dg\_wave}{rs}
\coderef{dg\_ccz4}{rs}
